\section{Background and Project Aim}
\subsection*{Background}
\
\begin{quote}
"When you don't give a f*** about the output of LLMs."
\end{quote}
\noindent
Large Language Models (LLMs) are becoming increasingly sophisticated, producing outputs that are not only more coherent but also more impactful in shaping digital interactions. As these models grow in size and capability, their influence expands \cite{hao2022languagemodelsgeneralpurposeinterfaces}.
\\
Generative AI, while impressive, operates by predicting the next word in a sequence based on patterns learned from large datasets. Resulting in no understanding of meaning of its outputs. This lack of comprehension and judgement becomes problematic when it generates toxic content. Often failing to recognise appropriate responses itself.
\\
As language is highly context dependent, so are the potentially toxic outputs (or model failures). When characterising these model failures we encounter two main challenges: benchmarking toxicity and explaining the model that generated the toxic output. 
In order to understand why a model generates toxic outputs there first has to be a definition of a toxic output. Therefore, benchmarking the toxicity in a certain piece of text is pivotal before diving deeper into the linguistics of the text. With the rise of LLMs, Explainable AI (XAI) is becoming increasingly important. Since toxicity in language can have a severe societal impact it is essential that the decision-making of these models are transparent. Thus, understanding the LLMs, which are mainly black box models, is a hard but important task.
\\
Given these challenges, LLMs should not simply be treated as tools where data is fed in and outputs are accepted without question. Their use must be accompanied by ethical considerations. Understanding how and why such toxicity arises is crucial. This context serves as the foundation of our project: \textbf{Characterising Toxicity in Language Models}, focusing on figuring out when and why some LLMs create toxic outputs.

\subsection*{Dataset}
Decoding Trust is a project aimed at providing a thorough assessment of trustworthiness in Generative Pretrained Transformer (GPT) models \cite{wang2024decodingtrustcomprehensiveassessmenttrustworthiness}. The project provides datasets to test models on a variety of aspects of trustworthiness, including toxicity. The prompts concerning toxicity are a subset of RealToxicityPrompts \cite{gehman2020realtoxicitypromptsevaluatingneuraltoxic}. The dataset consists of prompts, along with their toxicity score, calculated using PerspectiveAPI \cite{lees2022newgenerationperspectiveapi}, a widely used toxicity-scoring tool developed by Google Jigsaw. This data will be used in this project to analyse models, as well as to analyse the features in prompts that solicit toxic responses.

\subsection*{Model Selection}
Not all LLMs are created equal, there are many different architectures an LLM can be built with, which come with its own quirks. For this reason part of the project will look into a number of models to find out when various models generate toxic content, and investigate the level of toxicity generated by various models. This project will only look into models that can be run locally, are from different research groups, and are the newest and most advanced models readily available as of December 2024. The models that were selected for this project are: MistalAI's Mixtral 8x7B \cite{jiang2024mixtralexperts}, Meta's llama3 \cite{grattafiori2024llama3herdmodels}, Google's Gemma \cite{gemmateam2024gemmaopenmodelsbased}, and BigScience's BLOOM \cite{workshop2023bloom176bparameteropenaccessmultilingual}.  

%\subsection{Project Aim}

% \noindent
% The goal of the project is to figure out when and why some LLM’s create toxic outputs. 